% answers_short.tex —— 2013 年·全国大纲卷·理科 · 简短答案速查
% 由 tools/make_short_answers.py 自 answers.tex 的 \gksolhead/\gkanswer 提取，
% 并与 2013/2013_数学真题_大纲理_答案要点.md 逐题比对后生成，勿手改（改源 MD 后重跑）。

\begin{gknotice}
\textbf{使用说明}：本节只给最终答案便于快速核对，完整推导见解析版。答案为编写者独立推导，\textbf{非官方标准答案}。
\end{gknotice}

\gkpart{选择题}
\begin{center}\begin{tabular}{ccccccccccccc}
\toprule
题号 & 1 & 2 & 3 & 4 & 5 & 6 & 7 & 8 & 9 & 10 & 11 & 12 \\
\midrule
答案 & B & A & B & B & A & C & D & B & D & A & D & C \\
\bottomrule
\end{tabular}\end{center}

\gkpart{填空题}
\begin{itemize}[leftmargin=2.6em]
  \item[\textbf{13.}] $2\sqrt2$。
  \item[\textbf{14.}] $480$。
  \item[\textbf{15.}] $\left[\dfrac12,4\right]$。
  \item[\textbf{16.}] $16\pi$。
\end{itemize}

\gkpart{解答题}
\begin{itemize}[leftmargin=2.6em]
  \item[\textbf{17.}] $a_n=2n-1$ 或 $a_n=3$（$n\in\mathbf N^*$）。
  \item[\textbf{18.}] (1) $B=\dfrac{2\pi}3$；(2) $C=\dfrac{\pi}{12}$ 或 $C=\dfrac\pi4$
  \item[\textbf{19.}] (1) 证明见解析；(2) $\theta=\arccos\left(-\dfrac{\sqrt6}3\right)=\pi-\arccos\dfrac{\sqrt6}3$，约为 $144.7^\circ$
  \item[\textbf{20.}] (1) $\dfrac14$；(2) $P(X=0)=\dfrac18$，$P(X=1)=\dfrac58$，$P(X=2)=\dfrac14$，$E(X)=\dfrac98$
  \item[\textbf{21.}] (1) $a=1$，$b=2\sqrt2$；(2) 证明见解析
  \item[\textbf{22.}] (1) $\lambda_{\min}=\dfrac12$；(2) 证明见解析
\end{itemize}
